% Vietnamese translation for OI Public Library
% Source-PDF: 动态规划 DYNAMIC PROGRAMMING/初探数位dp.pdf
\documentclass[11pt]{article}
\usepackage{oipl-vn}

\title{Nhập môn quy hoạch động chữ số}
\author{Lazycal}
\date{}

\begin{document}
\maketitle

\origtitle{Khảo sát ban đầu về các bài toán thống kê theo chữ số}
\sourcepath{Dynamic Programming / Introductory digit DP.pdf}

\begin{abstract}
Tài liệu nguồn là một bộ slide ảnh về các bài toán đếm theo chữ số. Bản dịch
chuyển nội dung thành ghi chú liền mạch: trước hết thống nhất ký hiệu khoảng,
sau đó trình bày cách tiền xử lý trạng thái theo số chữ số và cách quét từng
chữ số của cận trên để tính số lượng trong một khoảng. Các slide bài tập và
mã giả được đối chiếu bằng OCR và kiểm tra trực quan.
\end{abstract}

\tableofcontents

\section{Kiến thức cơ bản}

Một số ký hiệu khoảng được dùng xuyên suốt tài liệu:
\[
  [l,r]=\{x\mid l\le x\le r\},\qquad
  [l,r)=\{x\mid l\le x<r\},
\]
\[
  (l,r]=\{x\mid l<x\le r\},\qquad
  (l,r)=\{x\mid l<x<r\}.
\]
Dấu ngoặc vuông nghĩa là lấy cả đầu mút; dấu ngoặc tròn nghĩa là không lấy
đầu mút đó.

Trong các kỳ thi OI, ta thường gặp yêu cầu đếm số lượng số nguyên thỏa mãn
một điều kiện nào đó trong đoạn $[l,r]$. Dạng này thường được đổi thành hiệu
của hai truy vấn tiền tố:
\[
  \operatorname{count}([l,r])
  =\operatorname{count}([0,r])-\operatorname{count}([0,l)).
\]
Vì vậy phần chính là xây dựng một hàm đếm các số trong $[0,n)$ hoặc
$[0,n]$.

\section{Tư tưởng và phương pháp cơ bản}

Với một số $x<n$, khi so sánh từ chữ số cao xuống chữ số thấp, luôn tồn tại
một vị trí đầu tiên tại đó chữ số của $x$ nhỏ hơn chữ số tương ứng của $n$.
Ví dụ với $n=58$ trong hệ thập phân:
\[
  x=49
\]
nhỏ hơn $n$ ngay ở hàng chục, còn
\[
  x=51
\]
giống $n$ ở hàng chục và nhỏ hơn $n$ ở hàng đơn vị.

Từ nhận xét này, ta quét các chữ số của $n$ từ cao xuống thấp, thử chọn chữ số
đầu tiên nhỏ hơn cận trên. Một khi vị trí đó đã được chọn, các chữ số phía sau
không còn bị ràng buộc bởi $n$ nữa; chúng có thể chạy từ $00\ldots0$ đến
$99\ldots9$ trong hệ thập phân. Vì phần đuôi đã tự do, ta tiền xử lý số cách
cho từng độ dài và từng trạng thái rồi cộng trực tiếp.

Khuôn mẫu trạng thái thường có dạng
\[
  f[i,\mathrm{st}],
\]
trong đó $i$ là số chữ số còn lại cần điền, còn $\mathrm{st}$ là trạng thái
phụ thuộc vào bài toán. Các số có thể có chữ số $0$ ở đầu; chẳng hạn
$00058$ vẫn được xem là một số có $5$ chữ số trong bảng tiền xử lý.

Nếu $i=4$, thì $f[i,\mathrm{st}]$ đếm các số từ $0000$ đến $9999$ thỏa điều
kiện tương ứng với trạng thái $\mathrm{st}$. Khi quyết định chữ số thứ $i$,
ta thử các giá trị có thể, cập nhật trạng thái mới $\mathrm{st}'$, rồi cộng:
\[
  f[i,\mathrm{st}] \mathrel{+}= f[i-1,\mathrm{st}'].
\]

\section{HDU 2089}

Đề bài yêu cầu đếm trong đoạn $[n,m]$ số lượng số không chứa chuỗi chữ số
\texttt{62} và cũng không chứa chữ số \texttt{4}. Ví dụ, $62315$ chứa
\texttt{62}, còn $88914$ chứa \texttt{4}, nên cả hai đều bị loại. Ràng buộc
trong slide là
\[
  0<n\le m<1000000.
\]

Theo phương pháp chung, ta tiền xử lý bảng $f$, sau đó trả lời bằng
\[
  \operatorname{solve}(m)-\operatorname{solve}(n-1),
\]
trong đó $\operatorname{solve}(u)$ đếm các số hợp lệ trong $[0,u]$.

\subsection{Tiền xử lý}

Đặt $f[i,j]$ là số lượng số có đúng $i$ chữ số, chữ số đầu là $j$, và không
chứa \texttt{62} hay \texttt{4}. Ví dụ $f[2,6]$ đếm các số hai chữ số bắt đầu
bằng $6$:
\[
  60,61,63,65,66,67,68,69.
\]
Hai lựa chọn $62$ và $64$ bị loại, lần lượt vì sinh chuỗi \texttt{62} và vì
có chữ số \texttt{4}.

Khởi tạo
\[
  f[0,0]=1.
\]
Với $i=1,\ldots,7$, thử chữ số đầu $j$ và chữ số kế tiếp $k$:
\[
  f[i,j]\mathrel{+}=f[i-1,k]
\]
nếu $j\ne 4$ và không đồng thời có $j=6,k=2$. Dạng vòng lặp là:
\begin{verbatim}
for i = 1 .. 7
  for j = 0 .. 9
    for k = 0 .. 9
      if j != 4 and not (j == 6 and k == 2)
        f[i][j] += f[i - 1][k]
\end{verbatim}

\subsection{Đếm tiền tố}

Giả sử các chữ số của $u$ được lưu từ phải sang trái:
\[
  \operatorname{digit}[1]\text{ là hàng đơn vị},\quad
  \operatorname{digit}[\mathrm{len}]\text{ là chữ số cao nhất}.
\]
Ví dụ $u=58$ thì $\operatorname{digit}[1]=8$ và
$\operatorname{digit}[2]=5$.

Khi quét từ $i=\mathrm{len}$ xuống $1$, ta thử mọi chữ số $j$ nhỏ hơn
$\operatorname{digit}[i]$ ở vị trí hiện tại. Nếu chọn $j$ hợp lệ với tiền tố
đã có, toàn bộ $i-1$ chữ số phía sau được đếm bằng bảng $f[i,j]$. Nếu chính
chữ số của $u$ đã tạo ra chữ số \texttt{4} hoặc chuỗi \texttt{62}, quá trình
dừng vì mọi số đi tiếp theo cùng tiền tố đều không hợp lệ.

\begin{verbatim}
for i = len .. 1
  for j = 0 .. digit[i] - 1
    if j != 4 and not (j == 2 and digit[i + 1] == 6)
      ans += f[i][j]
  if digit[i] == 4 or (digit[i] == 2 and digit[i + 1] == 6)
    break
\end{verbatim}

\section{HDU 3652}

Đề bài yêu cầu đếm các số nhỏ hơn $n$, chia hết cho $13$, và biểu diễn thập
phân có chứa chuỗi \texttt{13}.

Ta vẫn tiền xử lý trước, rồi quét theo chữ số để thống kê tiền tố. Vì bài toán
đòi hỏi đồng thời ``đã chứa \texttt{13} hay chưa'' và phần dư modulo $13$, ta
dùng trạng thái:
\[
  f[i,j,k,l].
\]
Ý nghĩa là: xét các số có $i$ chữ số, chữ số đầu là $j$; biến $k$ cho biết
đã chứa chuỗi \texttt{13} hay chưa; biến $l$ là phần dư modulo $13$.

Khi quyết định chữ số thứ $i$, thử $x=0,\ldots,9$. Nếu đã chứa \texttt{13}
thì phần đuôi có thể ở trạng thái đã chứa \texttt{13}; nếu chưa, cặp chữ số
$j=1,x=3$ sẽ chuyển trạng thái sang đã chứa. Phần dư được cập nhật bằng giá
trị của chữ số mới ở hàng $10^{i-1}$:
\[
  l' \equiv l-j\cdot 10^{i-1}\pmod{13}.
\]
Các chuyển trong slide có dạng:
\begin{align*}
  f[i,j,1,l] &\mathrel{+}= f[i-1,x,1,l'],\\
  f[i,j,1,l] &\mathrel{+}= f[i-1,x,0,l']
    \quad\text{nếu }j=1,\ x=3,\\
  f[i,j,0,l] &\mathrel{+}= f[i-1,x,0,l']
    \quad\text{nếu }(j,x)\ne(1,3).
\end{align*}

Khi đếm các số nhỏ hơn $n$, ta quét từng chữ số giống bài trước. Cần ghi nhớ
phần dư của tiền tố đã chọn và cờ cho biết tiền tố đã xuất hiện \texttt{13}
hay chưa. Với mỗi chữ số thử nhỏ hơn cận, bảng $f$ trả lời số lượng phần đuôi
khi phần dư cuối cùng phải bằng $0$ modulo $13$.

Một dạng mã của bước đếm trong slide:
\begin{verbatim}
bit[0] = 1;
for (ll i = 1; i <= 12; ++i) bit[i] = bit[i - 1] * 10;

for (ll i = digit[0], mod = 0; i; --i) {
  for (ll j = 0; j < digit[i]; ++j) {
    ans += f[i][j][1][(13 - mod * bit[i] % 13) % 13];
    if (t || (j == 3 && digit[i + 1] == 1))
      ans += f[i][j][0][(13 - mod * bit[i] % 13) % 13];
  }
  if (digit[i + 1] == 1 && digit[i] == 3) t = 1;
  mod = (mod * 10 + digit[i]) % 13;
}
\end{verbatim}

\section{URAL 1057}

Đề bài cho đoạn $[X,Y]$ và yêu cầu đếm số nguyên thỏa mãn: số đó bằng tổng của
đúng $K$ lũy thừa nguyên không âm khác nhau của $B$. Ví dụ với
$X=15,Y=20,K=2,B=2$, ba số thỏa mãn là
\[
  17=2^4+2^0,\qquad
  18=2^4+2^1,\qquad
  20=2^4+2^2.
\]
Ràng buộc:
\[
  1\le X\le Y\le 2^{31}-1,\qquad
  1\le K\le 20,\qquad
  2\le B\le 10.
\]

Một số là tổng của các lũy thừa khác nhau của $B$ khi và chỉ khi trong biểu
diễn cơ số $B$, mỗi chữ số chỉ là $0$ hoặc $1$, và có đúng $K$ chữ số bằng
$1$. Vì vậy ta có thể chuyển mọi cơ số $B$ về cùng một bài toán nhị phân trên
các chữ số $0,1$.

Đặt
\[
  \operatorname{cnt}[i,j]
\]
là số lượng xâu nhị phân dài $i$ có đúng $j$ bit bằng $1$. Tiền xử lý:
\[
  \operatorname{cnt}[i,j]
  =\operatorname{cnt}[i-1,j]+\operatorname{cnt}[i-1,j-1].
\]
Để tính $\operatorname{count}([0,n])$, đổi $n$ sang cơ số $B$ rồi quét từ trái
sang phải. Mỗi khi gặp chữ số $1$, ta có thể đặt vị trí hiện tại thành $0$ và
cộng số cách điền phần sau với đúng số bit $1$ còn thiếu:
\[
  \text{nếu }\operatorname{digit}[i]=1
  \text{ thì } ans\mathrel{+}=\operatorname{cnt}[i,\mathrm{need}].
\]
Ở đây $\mathrm{need}$ là số chữ số $1$ còn cần ở phần đuôi. Nếu gặp chữ số lớn
hơn $1$, ta chọn chữ số hiện tại là $1$ rồi đặt mọi chữ số phía sau tự do
trong $\{0,1\}$; sau đó không cần tiếp tục bám sát cận, vì tiền tố đã nhỏ hơn
$n$.

Vấn đề cuối cùng là xử lý cơ số khác $2$. Với một truy vấn $n$, ta cần tìm số
lớn nhất không vượt quá $n$ mà biểu diễn cơ số $B$ chỉ chứa $0$ và $1$. Cách
làm là tìm chữ số đầu tiên từ trái sang phải khác $0,1$, đổi nó thành $1$ và
đặt toàn bộ chữ số bên phải thành $1$. Biểu diễn cơ số $B$ thu được có thể
được xem như một xâu nhị phân để trả lời truy vấn. Ví dụ:
\[
  n=(10204)_9 \quad\Longrightarrow\quad n'=(10111)_2
\]
trong không gian đếm các chữ số hợp lệ.

\section{Một ví dụ với phép xor}

Bài \texttt{test-09-07-p1} trong slide phát biểu: cho dãy $A[i]$ độ dài $n$,
tính tổng các giá trị
\[
  A[i]\operatorname{xor}A[j]\qquad (i<j).
\]
Ý tưởng cũng là xử lý từng bit. Với mỗi bit, khi xét một phần tử, ta biết
trước đó có bao nhiêu phần tử có bit bằng $0$ và bao nhiêu phần tử có bit
bằng $1$; bit hiện tại đóng góp vào tổng đúng khi hai bit khác nhau.

\section{Mở rộng: SPOJ Sorted Bit Sequence}

Slide gợi ý bài \texttt{SORTBIT}:
\[
  \text{\url{http://www.spoj.pl/problems/SORTBIT}}
\]
hoặc bản trên HUST:
\[
  \text{\url{http://acm.hust.edu.cn/vjudge/problem/viewProblem.action?id=18852}}.
\]
Nội dung và phân tích của bài này được dẫn tới bài viết tham khảo trong phần
cuối.

\section{Kết luận}

Cốt lõi của các bài toán chữ số là tư tưởng ``xác định từng vị trí''. Vì số
lượng thao tác cơ bản chỉ ở mức $O(\log n)$ theo giá trị cận trên, ta có thể
chấp nhận trạng thái phụ phức tạp hơn một chút, hoặc dùng chia đôi, liệt kê
và các kỹ thuật con khác để giảm độ phức tạp tư duy và cài đặt.

Khi cần đếm trong đoạn $[l,r]$, dạng thường dùng vẫn là:
\[
  \operatorname{count}([l,r])
  =\operatorname{count}([0,r])-\operatorname{count}([0,l)).
\]

\section*{Tài liệu tham khảo}

\begin{itemize}
  \item Liu Cong, ``Bàn về các bài toán thống kê theo chữ số'', trong tuyển
  tập thuật toán được dẫn bởi slide.
  \item \url{http://hi.baidu.com/billdu/item/c749952ab2ab50c2ef10f137}
\end{itemize}

\end{document}
